\documentclass[version=preprint]{iacrcc}

\license{CC-by}

\title[running = {Equalis Committee Edition},
       subtitle = {Condensed Draft}]
      {Equalis: Deterministic Credit, Unified Liquidity, and Protocol-Perpetual Leverage}

\addauthor[orcid    = {0009-0009-7870-1337},
           inst      = {1},
           onclick   = {https://www.equalfilabs.com},
           email     = {mhooft@equalfilabs.com},
           surname   = {Hooft}
          ]{Matthew L. Hooft}

\addaffiliation[street = {}]{EqualFi Labs}

\begin{document}
\maketitle

\begin{abstract}
 Equalis offers: deterministic self-secured credit, explicit peer-to-peer credit (term and rolling), synthetic options via Direct agreements, unified per-token liquidity with encumbrance, dual-index incentives, EqualIndex fixed-composition baskets with basket flash loans, AMM auction venues, and Maker Auction Market (MAM) curves. The core design claim is narrow: protocol rules do not force collateral sale solely because of mark-to-market price movement. Continuation is governed by explicit obligations and deterministic state transitions. The architecture combines one canonical base pool per token, Position NFT accounting, module-level encumbrance namespaces, and shared fee routing into Fee Index and Active Credit Index rails. A key differentiator is \textbf{dual-sided ACI}: under protocol eligibility rules, ACI can reward both borrower-side deployed principal (including borrowed notional from 0\% self-secured credit) and lender-side deployed principal in P2P agreements (in addition to contractual lender interest). The paper does not claim risk elimination, guaranteed returns, or full MEV elimination; it claims legible risk, deterministic enforcement, reduced liquidity/incentive fragmentation, and elimination of reserve-repricing sandwich attacks on MAM curves.
\end{abstract}

\section*{Primitive Coverage at a Glance}
\begin{itemize}
  \item \textbf{Deterministic Core Credit:} same-asset self-secured borrowing at 0\% interest in protocol terms (with explicit action/maintenance fees in implementation), oracle-free solvency, no price-triggered liquidation.
  \item \textbf{Equalis Direct (P2P):} explicit bilateral credit with onchain offers, term and rolling agreements, optional lender-call covenant, deterministic default handling.
  \item \textbf{Synthetic Options:} physically settled option-like payoffs implemented as Direct agreements with repay-versus-forfeit decision logic.
  \item \textbf{Unified Liquidity:} one canonical base pool per token (plus optional managed pools) with module allocations represented via encumbrance rather than silo migration.
  \item \textbf{Dual Indexes:} per-token Fee Index (passive) and Active Credit Index (active) fed by shared fee routers across modules; ACI is designed to reward both active borrowing-side and active lending-side deployed capital.
  \item \textbf{EqualIndex:} static, fully backed baskets with deterministic mint/burn and proportional basket flash loans.
  \item \textbf{Tokenized Derivatives (ERC-1155):} standardized options and futures/forwards series with fungible per-series units, physical settlement, and deterministic reclaim windows.
  \item \textbf{Execution Surfaces:} solo/community AMM auctions and MAM time-based curves with deterministic quoting and explicit fee splits.
  \item \textbf{Protocol-Perpetual Leverage:} a protocol property, not a product type; no rule that forces unwind solely on price movement.
\end{itemize}

\section{Problem Framing}

Most DeFi credit is reactive: a health factor is continuously tested against external prices and violations are resolved through forced liquidation. This pattern can function, but it introduces:
\begin{itemize}
  \item operational burden (constant monitoring),
  \item correlated stress behavior (many liquidations at once),
  \item liquidation-path extraction during volatility,
  \item dependence on oracle governance, feed quality, and update timing.
\end{itemize}

At the same time, product expansion often fragments liquidity and incentives into independent silos. Capital moves between venues, fee rails split, and systems rely on emissions to patch short-term participation.

Equalis is designed around a different baseline: deterministic accounting for core solvency, explicit agreements for cross-asset risk, and unified per-token liquidity/fee rails across modules.

\section{System Architecture}

\subsection{One Canonical Pool Per Token}

For each token $T$, Equalis maintains a canonical base pool $\mathcal{P}_T$. Core features involving $T$ reference the same pool state and the same fee router for $T$.

\paragraph{Managed pools (implementation).}
The implementation also supports managed pools for the same underlying asset with custom controls and optional whitelist membership. Managed pools can route a configurable system share of fees back to the base rail. This preserves shared token-level fee plumbing while allowing curated pool variants.

\subsection{Position NFT as Account Container}

Each user position is represented by a Position NFT that tracks:
\begin{itemize}
  \item per-pool balances,
  \item same-asset debt,
  \item Direct agreement commitments,
  \item module encumbrances,
  \item index snapshots and claimable accrual.
\end{itemize}

Transferring the NFT transfers the balance sheet, subject to constraints required by active obligations and solvency rules.

\subsection{Encumbrance Instead of Liquidity Exit}

Modules consume inventory through encumbrance accounting:
\[
B^{free}_{p,T} \rightarrow B^{enc}_{p,T}
\]
under explicit module namespaces. Encumbered units are unavailable for conflicting actions until released. This enables module composability without duplicating token pools.

\subsection{Shared Fee Router and Dual Index Rails}

For each token $T$, fee events can route to:
\begin{itemize}
  \item direct participant payouts,
  \item treasury allocation,
  \item Fee Index rail ($I_T$),
  \item Active Credit Index rail ($A_T$).
\end{itemize}

Index updates follow:
\[
I_T \leftarrow I_T + \frac{F^{FI}_T}{W^{FI}_T}, \qquad
A_T \leftarrow A_T + \frac{F^{ACI}_T}{W^{ACI}_T}.
\]

This is the core anti-fragmentation mechanism: multiple modules can produce fees while feeding the same per-token rails.

\section{Deterministic Core Credit}

\subsection{Self-Secured Borrowing}

For token $T$, same-asset borrowing is bounded by:
\[
D_{p,T} \le \lambda_T \cdot \text{Deposit}_{p,T},
\]
with collateral and debt both denominated in $T$. Because solvency is same-denomination, core enforcement does not require external price conversion.

\subsection{0\% Interest in Protocol Terms}

Self-secured credit is specified as 0\% interest in the protocol model. Implementation still applies explicit deterministic fees (for example action and maintenance flows), but not oracle-driven variable interest tied to market pricing.

\subsection{No Price-Triggered Liquidation in Core}

The deterministic claim is specific:
\begin{itemize}
  \item no oracle threshold that auto-sells collateral solely on price moves,
  \item no liquidation auction path as default solvency mechanism for self-secured credit,
  \item continuation governed by explicit rule compliance in same-asset terms.
\end{itemize}

\section{Equalis Direct: Explicit P2P Credit}

\subsection{Two Agreement Shapes}

Equalis Direct exposes two agreement shapes for cross-asset credit:
\begin{itemize}
  \item \textbf{Term agreements:} fixed due timestamp with explicit grace/default logic.
  \item \textbf{Rolling agreements:} recurring payment interval with moving due timestamp (\texttt{nextDue}); schedule advancement is contingent on bringing obligations current.
\end{itemize}

This is an important distinction: ``perpetual'' is not a separate agreement type; persistence is a protocol behavior of compliant rolling/contractual credit, not a third loan class.

\subsection{Offer-Driven Onchain Discovery}

Both sides can post offers. The protocol supports lender- and borrower-posted structures, including tranche and ratio-tranche flows for repeated fills. Offers and lifecycle transitions are emitted as indexable events.

This is intended to support \textbf{CLOB-like discovery} for credit:
\begin{itemize}
  \item outstanding offers can be indexed by asset pair, collateral ratio, cadence, and covenant flags,
  \item fills, cancellations, and state transitions create a public event stream for order-book-style views,
  \item matching and settlement remain deterministic onchain even when discovery and ranking are rendered offchain.
\end{itemize}

\subsection{Explicit Covenants}

Optional covenants are agreement fields, including lender callability (\texttt{allowLenderCall}). If callability is not accepted in terms, lenders cannot unilaterally accelerate.

\subsection{Default as Contractual Event}

Defaults are handled as deterministic agreement transitions (obligation + grace based), not as oracle-liquidation competitions. Consequences are explicit in collateral routing and fee split rules.

\section{Synthetic Options via Direct Agreements}

Equalis Direct can express option-like payoffs without cash-settlement oracles by mapping terminal borrower choice to exercise logic:
\begin{itemize}
  \item \textbf{Repay:} return borrowed asset, recover collateral.
  \item \textbf{Forfeit/Exercise:} keep borrowed asset, forfeit collateral.
\end{itemize}

\subsection{Strike Encoding}

For borrowed asset $B$ and collateral asset $C$:
\[
K_{\text{implied}} = \frac{\texttt{collateralLockAmount}}{\texttt{principal}}.
\]
Strike is ratio-encoded via fixed quantities, not continuously oracle-priced.

\subsection{Premium and Style}

Premium can be represented as upfront interest/premium at origination. Agreement flags encode early exercise and early repay behavior, allowing American-like or European-like constraints.

\subsection{Settlement Model}

Settlement is physical. On exercise/default, collateral routing is deterministic and can split between lender compensation and protocol/index rails.

\subsection{Options Discovery and Event Surfaces}

Synthetic options inherit Direct's offer/event flow, so strike-like ratios, premium terms, expiry windows, and lifecycle outcomes are indexable from emitted events. This enables \textbf{CLOB-like options discovery} without requiring centralized matching logic: indexers and UIs can construct books of option-like intents while execution remains agreement-native and deterministic.

\subsection{ERC-1155 Tokenized Options and Futures/Forwards (Implementation)}

In addition to agreement-native Direct instruments, the implementation includes standardized ERC-1155 series modules for derivatives. The design goal is to support fungible lots and cleaner secondary transfer when users want product standardization instead of bespoke bilateral terms.

\paragraph{Token model.}
\texttt{OptionToken} and \texttt{FuturesToken} are ERC-1155 contracts where each derivative series corresponds to a token id. Units are fungible within the id and transferable with normal ERC-1155 semantics, while mint/burn authority remains protocol-controlled through the derivatives facets so lifecycle transitions remain deterministic.

\paragraph{Option series lifecycle.}
An options series is instantiated with explicit terms (underlying/quote pools, strike, expiry, style flags, and size). Creation locks maker collateral and mints option units for that series id. Exercise burns option units and settles physically according to series terms. After expiry plus settlement windows, unexercised residual collateral is reclaimable by the maker.

\paragraph{Futures/forwards series lifecycle.}
A futures/forwards series is instantiated with explicit forward price, expiry, settlement settings, and size. Creation locks maker inventory and mints ERC-1155 futures units. Settlement burns futures units and enforces deterministic delivery-versus-payment behavior. Maker reclaim is available only after the configured post-expiry grace logic.

\paragraph{Relationship to Direct agreements.}
Direct synthetic agreements are optimized for bespoke, covenant-rich bilateral contracts and Position-level portfolio transfer. ERC-1155 derivatives are optimized for standardized series, fungibility, and transferability per series id. Both surfaces are physically settled and feed the same unified pool/encumbrance accounting model.

\section{Unified Liquidity and Encumbrance Model}

\subsection{Why This Matters}

In fragmented systems, each new venue requires new liquidity and separate incentives. In Equalis, modules allocate from shared pools via encumbrance. This yields:
\begin{itemize}
  \item one token liquidity substrate,
  \item coherent accounting across modules,
  \item fee routing coherence across product surfaces.
\end{itemize}

\subsection{Encumbrance Safety Property}

Non-double-use must hold:
\begin{itemize}
  \item encumbered balances cannot be withdrawn as free liquidity,
  \item encumbered balances cannot simultaneously satisfy conflicting module states.
\end{itemize}

\subsection{Expansion Models via Encumbrance (Implementation Roadmap)}

The encumbrance system allows us to add new products without creating new liquidity silos. The key architectural point is that encumbrance is not just a safety primitive; it is an \textbf{expansion primitive} for unified rails.

\paragraph{Concrete expansion models.}
\begin{itemize}
  \item \textbf {Isolated Lending Markets:} lender allocations and borrower collateral are isolated per-market while still using pool-native principal and encumbrance buckets.
  \item \textbf{Isolated Perpetual Pools:} LP underwriting and trader margin are isolated through \texttt{directLent}/\texttt{directLocked}, with per-market PnL containment and shared pool accounting.
  \item \textbf{Isolated Prediction Markets:} LP capacity and trader stake are encumbered in bounded market state machines, then released at settlement, while fee routing remains connected to token rails.
\end{itemize}

\paragraph{Total-capture implication.}
Because modules allocate through the same per-token encumbrance and settle through the same fee router, the system has a path to \textbf{full economic capture} of product activity: new products can plug into the same Treasury/Fee Index/ACI rails instead of leaking value into isolated incentive islands. In this sense, ``add any product'' means: if a module can (i) lock/release inventory via encumbrance, (ii) enforce deterministic state transitions, and (iii) route fees through canonical per-token routers, it can be composed into the same value rails.

\section{Dual Index Incentive System}

\subsection{Fee Index (Passive)}

A conservative passive weight form is:
\[
w^{FI}_{p,T} = \max(\text{Deposit}_{p,T} - \text{Debt}_{p,T}, 0),
\]
to reduce borrow-deposit circular overcounting.

\subsection{Active Credit Index (Active)}

ACI weights track eligible deployed capital categories with maturity/anti-cycling controls in implementation. The intended economic split is explicitly dual-sided:
\begin{itemize}
  \item \textbf{Borrower-side ACI:} eligible borrowed principal can contribute active weight, including borrowed notional from self-secured 0\% credit when used under protocol-defined active categories.
  \item \textbf{Lender-side ACI:} eligible lender principal deployed into Direct agreements can also contribute active weight.
\end{itemize}

This means a P2P lender can earn contractual interest from the agreement \emph{and} earn ACI exposure on eligible lent capital through the shared index rail.

\subsection{Anti-Gaming Intent}

The dual-rail split is meant to keep passive and active incentives legible and harder to game via trivial cycling.

\subsection{Time Weighting and Self-Secured Loop Resistance}

ACI uses a \textbf{hard maturity gate} with weighted re-aging on increases. In the current implementation, active weight for a tracked state (debt-side or encumbrance-side) is:
\[
w^{ACI}_{p,T}(t)=
\begin{cases}
0, & \text{if age} < 24\text{h}\\
\text{principal}, & \text{if age} \ge 24\text{h}
\end{cases}
\]
When principal is increased, the state is re-aged by weighted dilution (not reset to a free full-weight state):
\[
\text{age}_{new}\approx
\frac{\text{principal}_{old}\cdot \text{age}_{old}}
{\text{principal}_{old}+\text{addedPrincipal}},
\]
with age capped by the gate horizon.

The anti-loop objective is explicit:
\begin{itemize}
  \item \textbf{Stop instant farming:} rapid borrow/repay or borrow/redeposit churn remains ineligible until maturity.
  \item \textbf{Neutralize trivial self-secured loops on passive rail:} Fee Index uses net passive weight (e.g., $\max(\text{Deposit}-\text{Debt},0)$), so borrowing against self and redepositing does not create free passive double-counting.
  \item \textbf{Constrain leverage recursion at solvency layer:} self-secured draws are LTV-bounded against net equity in deterministic same-asset accounting.
  \item \textbf{Reward sustained deployment:} ACI favors capital that remains actively deployed over time (borrowed or lent), not momentary notional spikes.
\end{itemize}

In short, loop resistance is layered: deterministic net-equity solvency limits draw capacity, Fee Index netting removes passive double-counting, and ACI maturity gating blocks immediate active-reward extraction from short-cycle churn.

\section{EqualIndex: Static Baskets and Basket Flash Loans}

\subsection{Primitive}

EqualIndex defines:
\begin{itemize}
  \item fixed underlying set $\mathcal{U}=\{T_1,\dots,T_k\}$,
  \item fixed per-unit quantities $\mathbf{q}=(q_1,\dots,q_k)$,
  \item no rebalance and no oracle NAV dependency.
\end{itemize}

\subsection{Mint/Burn Determinism}

Mint $x$ units deposits $x \cdot q_i$ of each $T_i$. Burn $x$ redeems pro-rata vault and fee-pot shares, less configured fees.

\subsection{Basket Flash Loans}

Borrower can atomically borrow proportional basket inventory:
\[
\text{FlashBorrow}(T_i) = x\cdot q_i, \qquad
\text{FlashRepay}(T_i) = x\cdot q_i + \phi_i(x).
\]
If repayment is incomplete in-transaction, operation reverts.

\subsection{Fee Split Coupling}

Per-underlying fees split between:
\begin{itemize}
  \item system share (treasury + index rails via underlying pool router),
  \item index-holder fee pots.
\end{itemize}

This explicitly couples basket activity to system rails rather than isolating it.

\section{Execution Surfaces}

\subsection{AMM Auctions}

Equalis includes solo and community auction modes where inventory comes from encumbered balances and fees split between makers and protocol rails.

\subsection{MAM Curves}

MAM provides one-sided, time-based pricing curves with deterministic linear interpolation:
\[
P(t)=
\begin{cases}
P_0, & t \le t_0 \\
P_1, & t \ge t_0 + D \\
P_0 + (P_1-P_0)\frac{t-t_0}{D}, & \text{otherwise}
\end{cases}
\]
supporting descending, ascending, and flat (limit-like) curves.

\subsection{MEV Posture}

MAM curves eliminate reserve-repricing (slippage-driven) sandwich attacks on the curve itself because fills do not move price as a function of trade size. MAM does not claim full MEV elimination; ordering competition, priority racing, and backrunning can still exist in a public mempool.

\section{Protocol-Perpetual Leverage}

\subsection{Definition}

\begin{quote}
A leveraged position is protocol-perpetual if protocol rules contain no transition that forces collateral sale solely due to external price movement.
\end{quote}

Termination is due to voluntary unwind or violation of explicit obligations (including opted-in callability where applicable), not automatic mark-to-market liquidation rules.

\subsection{Protocol Property Versus Product Type}

``Perpetual'' is not a third loan product. Equalis Direct exposes term and rolling agreements. Protocol-perpetual leverage describes what the rule set does \emph{not} do: it does not contain a price-triggered auto-close path. In practice:
\begin{itemize}
  \item term agreements can still be used in chained strategies, but each agreement has explicit maturity,
  \item rolling agreements are the natural vehicle for continuity because obligations are periodic and state can continue while compliant,
  \item continuation is a contractual state outcome, not an oracle-threshold state outcome.
\end{itemize}

\subsection{Why Rolling Agreements Enable Continuity}

Rolling agreements track contractual time via a moving due timestamp (\texttt{nextDue}) and an explicit grace window. The continuity logic is obligation-driven:
\begin{itemize}
  \item when required payments are made and obligations are brought current, schedule state advances,
  \item if obligations are missed beyond grace, default recovery paths become callable,
  \item at no point is ``price moved'' itself a sufficient condition for forced collateral sale by protocol rule.
\end{itemize}

This changes the source of fragility in leveraged strategies. Fragility shifts from mark-to-market oracle events to payment-discipline and counterparty-contract risk.

\subsection{Canonical Loop Construction}

Let:
\begin{itemize}
  \item $T_1$ be the target exposure asset,
  \item $T_2$ be the borrowed asset,
  \item $q_k$ be borrowed $T_2$ per loop iteration $k$,
  \item $\pi_k$ be $T_1$ received from execution in iteration $k$ net of execution costs.
\end{itemize}

One representative loop:
\begin{enumerate}
  \item borrow $q_k$ via a Direct rolling agreement in $T_2$,
  \item swap $q_k$ into $T_1$ via internal execution (for example AMM auction or MAM route),
  \item redeposit $\pi_k$ into $\mathcal{P}_{T_1}$,
  \item repeat while obligations are met and counterparties keep providing terms.
\end{enumerate}

After $N$ iterations, gross position-side exposure growth in $T_1$ is:
\[
\Delta E_{T_1}(N)=\sum_{k=1}^{N}\pi_k,
\]
while contractual liability stack is the sum of active borrowed principal and accrued contractual obligations in $T_2$ terms.

Under the ACI design, the borrower-side borrowed notional in compliant rolling strategies can also be eligible active weight, so loop economics are not only directional exposure plus fees, but also potential ACI accrual on active borrowed deployment.

However, this is \emph{not} intended to pay instantaneous loop churn: time-weighting and maturity gating are designed to suppress short-lived self-secured cycling and prioritize sustained active deployment.

\subsection{Operational State Machine}

From a protocol perspective, the loop can be described as:
\begin{enumerate}
  \item \textbf{Active-compliant:} obligations current, strategy continues.
  \item \textbf{Active-delinquent:} obligations not current, grace window still open.
  \item \textbf{Recovered/defaulted:} grace exceeded or explicit trigger reached; deterministic recovery paths execute.
  \item \textbf{Voluntary unwind:} borrower closes by repayment/unwind.
\end{enumerate}

Crucially, no additional state is triggered solely by external mark-to-market movement.

\subsection{Termination Conditions}

A protocol-perpetual loop persists until one of the following occurs:
\begin{itemize}
  \item \textbf{Voluntary unwind:} borrower reduces or closes leverage by repayment.
  \item \textbf{Default:} borrower misses explicit obligations beyond grace.
  \item \textbf{Opted-in callability:} lender accelerates only where callability was accepted in contract terms.
  \item \textbf{Strategy choice:} borrower stops re-looping because terms/execution are unattractive.
\end{itemize}

Notably absent is the common liquidation-system terminal condition:
\begin{quote}
``Price crossed threshold, therefore protocol forcibly sells collateral.''
\end{quote}

\subsection{Feasibility Versus Forced Unwind}

Protocol-perpetual does not mean frictionless. Market and credit conditions still matter:
\begin{itemize}
  \item execution quality (\texttt{minOut}, liquidity depth, spreads, fees),
  \item availability and pricing of rolling offers,
  \item borrower cashflow discipline for recurring obligations,
  \item counterparty appetite through volatility.
\end{itemize}

These factors determine whether continuation is attractive or possible, but they are distinct from protocol-level forced unwind rules conditioned solely on price.

\subsection{Worked Example (Qualitative)}

Suppose a borrower starts with $2$ units of $T_1$ and repeatedly borrows fixed-size tranches of $T_2$ through rolling agreements:
\begin{enumerate}
  \item each tranche is swapped into $T_1$ and redeposited,
  \item deposited base grows, improving strategy capacity for future iterations,
  \item periodic rolling obligations must be serviced to keep agreements current.
\end{enumerate}

If market price of $T_1$ drops, the protocol does not auto-liquidate the position based on that mark alone. The loop can continue through drawdown \emph{if} contractual payments continue and counterparties maintain offered terms. If payments fail beyond grace, deterministic recovery rules apply regardless of whether mark-to-market later recovers.

\subsection{Capital Efficiency Link}

The loop is system-relevant because it composes with unified pools and shared fee rails:
\begin{itemize}
  \item borrowed capital drives execution volume,
  \item execution and agreement fees route back into token-level rails,
  \item active deployment is represented in explicit accounting rather than opaque leverage multipliers.
\end{itemize}

This makes the loop an accounting- and fee-coupled behavior inside the architecture, not an isolated strategy external to incentive plumbing.

\subsection{Risk Accounting for the Loop}

Protocol-perpetual leverage still carries major risk surfaces:
\begin{itemize}
  \item \textbf{Counterparty risk:} lenders may tighten terms or refuse rollover.
  \item \textbf{Cashflow risk:} borrower may fail scheduled obligations.
  \item \textbf{Execution risk:} slippage and fee drag can dominate expected edge.
  \item \textbf{Smart contract risk:} invariants can fail if implementation is faulty.
  \item \textbf{Liquidity risk:} thin markets can block practical unwind routes.
\end{itemize}

Therefore, ``perpetual'' should be read as a statement about liquidation rule structure, not about guaranteed survival, guaranteed returns, or immunity to adverse conditions.

\subsection{What This Does Not Claim}

\begin{itemize}
  \item no guarantee of profitability,
  \item no guarantee of favorable execution conditions,
  \item no guarantee that counterparties keep offering acceptable terms,
  \item no elimination of default, smart contract, or market risk.
\end{itemize}

\section{Risk and Threat Model}

Equalis explicitly assumes adversarial users, adversarial ordering, and adversarial counterparties. Main risk classes:
\begin{itemize}
  \item smart contract and upgrade risk,
  \item token integration risk (non-standard ERC20 behavior),
  \item counterparty default risk in P2P credit,
  \item execution and liquidity risk in market modules,
  \item MEV and ordering risk,
  \item accounting/index manipulation attempts.
\end{itemize}

Mitigation strategy is deterministic enforcement, explicit contracts, accounting isolation, and module-scoped blast-radius controls.

\section{Governance Path: Early Mutability to Progressive Freezing}

Equalis treats upgradeability as temporary scaffolding. The intended sequence:
\begin{enumerate}
  \item experimental iteration and hardening,
  \item stabilization and audit depth,
  \item selective freezing of mature surfaces,
  \item minimized bounded knobs and walkaway-oriented operation.
\end{enumerate}

The implementation uses Diamond/facet modularity and owner/timelock governance machinery during early stages, with a target of reducing discretionary control over time.

\section{Implementation Mapping and Invariants}

\subsection{Current Architecture (Implementation)}

Reference implementation is Diamond-based with facets spanning:
\begin{itemize}
  \item pools, positions, and deterministic credit,
  \item Direct agreements and rolling lifecycle/payment flows,
  \item ERC-1155 tokenized options and futures/forwards series modules,
  \item execution modules (AMM auctions, MAM),
  \item EqualIndex operations,
  \item admin/timelock configuration,
  \item view/preview read paths.
\end{itemize}

\subsection{Core Invariants}

\begin{itemize}
  \item \textbf{Conservation:} pool totals reconcile with user balances and explicit fee buffers.
  \item \textbf{Non-double-use:} encumbered units cannot satisfy conflicting states.
  \item \textbf{Deterministic bound safety:} self-secured debt respects configured same-asset constraints.
  \item \textbf{Index monotonicity:} $I_T$ and $A_T$ are non-decreasing and fee-event-driven.
  \item \textbf{Agreement determinism:} lifecycle transitions follow explicit permitted edges.
  \item \textbf{EqualIndex backing:} supply maps to backing plus explicit fee-pot accounting.
\end{itemize}

\section{Claims and Non-Claims (Compact)}

\subsection*{Claims}
\begin{itemize}
  \item deterministic core credit without oracle liquidations,
  \item explicit legible cross-asset risk through bilateral contracts,
  \item unified liquidity and unified fee rails per token,
  \item dual-sided ACI where both eligible borrowed and lent capital can earn active index rewards,
  \item composable module expansion without mandatory liquidity silo duplication,
  \item protocol-perpetual leverage as a rule-level property.
\end{itemize}

\subsection*{Non-Claims}
\begin{itemize}
  \item no promised yield or profitability,
  \item no elimination of smart contract/counterparty risk,
  \item no full MEV elimination guarantee beyond reserve-repricing sandwich attacks on MAM curves,
  \item no claim that strategy feasibility is price-independent.
\end{itemize}

\section{Conclusion}

Equalis proposes a deterministic architecture where core solvency is not delegated to oracle-triggered liquidation, cross-asset risk is contractual and explicit, and module expansion feeds shared per-token fee rails instead of fragmenting them. A central differentiator is that active rewards are not single-sided: the same ACI rail is designed to reward eligible borrowing-side deployment and eligible lending-side deployment, including 0\% self-secured borrowed notional and Direct lender principal under protocol rules. The practical contribution is not a claim of risk removal; it is a stricter separation between protocol rule determinism and market risk, with enforcement legibility and a path toward progressively immutable operation.

\paragraph*{Note on Versioning}
This document is a condensed committee edition. The full version is available at: https://github.com/EqualFiLabs/EqualFi

\end{document}
